% Text of the poster. The figures come from figures/, made by poster_figures.py.
\newcommand{\TitleFirst}{Completing}
\newcommand{\TitleSecond}{broken meshes}
\newcommand{\Author}{Lucas Allara}
\newcommand{\Subtitle}{An implementation of Jacobson, Kavan and Sorkine-Hornung (SIGGRAPH 2013),\\
  extended with cap fairing and symmetry completion}
\newcommand{\Course}{Computer Graphics}
\newcommand{\Date}{November 2026}
\newcommand{\ColInput}{Input}
\newcommand{\ColMethod}{Method of [1]}
\newcommand{\ColFaired}{With fairing}
\newcommand{\ColSymmetry}{With fairing and symmetry}
\newcommand{\ColSymmetryShort}{With symmetry}
\newcommand{\ColOriginal}{Original}

\newcommand{\SecIntro}{Introduction}
\newcommand{\SecContribution}{Contribution}
\newcommand{\SecResults}{Results}
\newcommand{\SecLimitations}{Limitations}
\newcommand{\SecConclusions}{Conclusions}
\newcommand{\SecReferences}{References}

% ---------------------------------------------------------------- introduction
\newcommand{\IntroTextA}{%
  Many real meshes have holes, overlapping pieces or loose faces, and simulating or printing
  a model requires knowing which points are inside. Jacobson et al.\ [1] use the generalized
  winding number, which adds up the signed solid angles that the oriented faces subtend as seen
  from \(p\).\par}
\newcommand{\Formula}{\(w(p) = \dfrac{1}{4\pi} \displaystyle\sum_{f} \Omega_f(p)\)}
\newcommand{\WLabelA}{closed mesh}
\newcommand{\WLabelB}{with a hole}
\newcommand{\WText}{%
  On a closed mesh, \(w\) is 1 inside and 0 outside. With a hole it varies smoothly between
  the two, and the level 0.5, dashed in the figure, already closes the sphere.\par}
\newcommand{\IntroTextB}{%
  The convex hull is tetrahedralized respecting the input faces [2]. A graph cut labels each
  tetrahedron inside or outside by its \(w\), keeping neighbors of similar \(w\) together.
  Two heuristics push the output, the boundary of the inside region, to be a manifold that
  contains the input faces.\par}

% ---------------------------------------------------------------- contribution
\newcommand{\ContributionText}{%
  Where surface was missing, the output of [1] is made of faces of the tetrahedralization,
  which close each hole with a cap of flat triangles and no detail. Two stages are added after
  the method, without modifying it.\par}
\newcommand{\FairingText}{%
  \textcolor{fairedtext}{\bfseries Fairing.} The input stays fixed. Each cap is subdivided
  and its vertices solve \(\Delta^3 x = 0\) [4], which continues the surface around it in
  position, tangent and curvature. A cap that overshoots is solved again with a lower order.
  TetWild [3] can smooth such caps with the Laplacian alone.\par}
\newcommand{\SymmetryText}{%
  \textcolor{symmetry}{\bfseries Symmetry.} The plane that best maps the input onto itself is
  found, ignoring the 20\% of the points that match worst so that the holes do not count. The
  mirrored faces that fall in the holes are added to the input [5] and the method runs
  again.\par}

% ---------------------------------------------------------------- results
\newcommand{\NumA}{−42\%}
\newcommand{\NumAText}{error in the holes with fairing, on average over the 18 cases.}
\newcommand{\NumB}{−78\%}
\newcommand{\NumBText}{error in the holes with fairing and symmetry, on average over the 18 cases.}
\newcommand{\NumC}{18/18}
\newcommand{\NumCText}{cases in which each stage lowers the error of the one before.}
\newcommand{\TeaserCaption}{Homer with 11\% of its surface removed. The method of [1] closes the holes with flat faces (error 2.07\%), fairing smooths the caps (1.51\%), and symmetry brings back the mouth and the hand (0.37\%).}
\newcommand{\RendersCaption}{Spot and Homer at the 22\% level, with 28.8\% and 23.1\% of the surface
  removed. Symmetry copies the hand and the foot of Homer from the other side. Its face is missing on
  both sides of the plane and stays smooth.}
\newcommand{\TableHead}{mean error in the holes}
\newcommand{\TableSpot}{\tabhead{Spot}, 5\,856 triangles}
\newcommand{\TableHomer}{\tabhead{Homer}, 12\,000 triangles}
\newcommand{\TableChange}{change}
\newcommand{\RowLaplacian}{fairing with \(\Delta\)}
\newcommand{\RowBilaplacian}{fairing with \(\Delta^2\)}
\newcommand{\RowFaired}{with fairing, \(\Delta^3\)}
\newcommand{\TableCaption}{The means of the chart and of two lower orders of the fairing, with the mean
  change over the 18 cases. With \(\Delta\), which only continues the position, the error does not go
  down, and \(\Delta^3\) is lower than \(\Delta^2\) in 13 of them.}
\newcommand{\ChartCaption}{Mean distance from the removed surface to the completed mesh, in percent of
  the longest side of the bounding box. Each point averages three random sets of five holes, and its
  bar spans them. The dotted line is the original measured against itself. The measure does not
  penalize extra surface.}

% ---------------------------------------------------------------- limitations and conclusions
\newcommand{\LimitationsText}{%
  Symmetry requires a symmetric model whose other side is intact. If surface is missing on
  both sides of the plane, as with the face of Homer at 22\%, there is nothing to copy from.
  Without symmetry, fairing completes the overall shape but not the small details.\par
  \vspace{2mm}
  A piece that a hole separates from the rest stays separated, because the cut does not
  control the topology. The mirrored faces are not stitched to the input. With them, the
  output of the second pass has loose pieces in 5 of the 18 cases and 4 non-manifold edges in
  2 cases.\par}
\newcommand{\ConclusionsText}{%
  The method of [1] closes the holes, but it does not give shape to what is missing. On Spot
  and Homer, a local fairing of the caps, which does not move the input, lowers the error in the
  holes by 42\% on average, as long as it continues the tangent, since with the Laplacian the
  error does not go down. Symmetry takes the reduction to 78\% by copying details that fairing
  cannot invent.\par
  \vspace{2mm}
  The implementation, in C++, completes a model in one to two seconds on 16 cores. Future work
  includes symmetries of parts of the model and a graph cut that controls the topology.\par}

% ---------------------------------------------------------------- references
\newcommand{\ReferencesText}{%
  \bib{1}A.~Jacobson, L.~Kavan, O.~Sorkine-Hornung. Robust inside-outside segmentation using
  generalized winding numbers. \emph{ACM TOG} 32(4), 2013.\par
  \bib{2}H.~Si. TetGen, a Delaunay-based quality tetrahedral mesh generator. \emph{ACM TOMS}
  41(2), 2015.\par
  \bib{3}Y.~Hu et al. Tetrahedral meshing in the wild. \emph{ACM TOG} 37(4), 2018.\par
  \bib{4}M.~Botsch, L.~Kobbelt. An intuitive framework for real-time freeform modeling.
  \emph{ACM TOG} 23(3), 2004.\par
  \bib{5}X.~Li et al. Symmetry and template guided completion of damaged skulls.
  \emph{Computers \& Graphics} 35(4), 2011.\par}
\newcommand{\Credits}{Spot is by K.~Crane and Homer comes from common-3d-test-models.
  Images and measurements by the author.}
